\chapter*{结束语：把每一次按键变成证据}
\addcontentsline{toc}{chapter}{结束语：把每一次按键变成证据}

输入法工程最迷人的地方，是抽象算法与具体系统行为在毫秒尺度相遇。一个候选分数可能来自百万级词库和图搜索，一个按键却仍可能因为事件路径不同而没有进入控制器。优秀的输入法开发者既能推导评分函数，也愿意在第二台机器的 Notes 里逐字验证。

MyIME 不是商业输入法能力的终点，却是很好的学习起点。读者不必照抄所有实现。真正需要带走的是一套方法：分离核心与系统、用最小模型解释状态、用可重复测试约束算法、用跨机器验收约束发布，并始终区分“设计意图”和“实测证据”。

\begin{summarybox}
如果你已经能够独立回答“这个按键在哪里被消费”“这个候选为什么排第一”“这个用户数据写到哪里”“这个 DMG 为什么在另一台机器可用”，那么你已经不只是完成了一个输入法，而是完成了一次小型系统工程训练。
\end{summarybox}

\chapter*{版本与勘误记录}
\addcontentsline{toc}{chapter}{版本与勘误记录}

本教材首版对应 MyIME 1.0.10。建议课程维护者在项目升级时记录：源码提交号、教材章节、受影响的代码清单范围、行为变化和重新验证结果。不要直接覆盖旧结论；教学材料需要保留“为什么过去这样写、后来为什么改变”的工程历史。

问题和勘误可通过项目仓库提交。报告应至少包含 macOS 版本、目标应用、输入序列、预期结果、实际结果和最小复现步骤。
